% Preámbulo compartido para los anexos publicados de forma independiente (Drive).
% Reproduce los paquetes y definiciones relevantes de main.tex para que cada
% anexo compile como documento autónomo, resolviendo las referencias cruzadas
% al documento principal mediante el paquete xr.

\usepackage[utf8]{inputenc}
\DeclareUnicodeCharacter{2715}{\texttimes}
\DeclareUnicodeCharacter{2192}{\ensuremath{\rightarrow}}
\DeclareUnicodeCharacter{26A0}{\ensuremath{\triangle}}
\usepackage[spanish,es-noshorthands]{babel}

\usepackage{float}
\usepackage{booktabs}
\usepackage{longtable}
\usepackage{array}
\usepackage{tabularx}
\usepackage{graphicx}
\usepackage{subcaption}
\usepackage{rotating}
\usepackage{colortbl}

\usepackage{url}

\usepackage{xcolor}
\definecolor{naranja}{rgb}{1,0.5,0}
\definecolor{rojo}{rgb}{1,0,0}
\definecolor{SteelBlue}{rgb}{0.3,0.5,0.7}
\definecolor{dkgreen}{rgb}{0,0.6,0}
\definecolor{gray}{rgb}{0.5,0.5,0.5}
\definecolor{mauve}{rgb}{0.58,0,0.82}
\definecolor{darkblue}{rgb}{0.0,0.0,0.6}
\definecolor{cyan}{rgb}{0.0,0.6,0.6}
\definecolor{claregreen}{RGB}{4,180,95}

\usepackage{listingsutf8}
\lstset{
  basicstyle=\footnotesize,
  numberstyle=\footnotesize,
  numbersep=4pt,
  backgroundcolor=\color{white},
  breaklines=true,
  breakatwhitespace=false,
  title=\lstname,
  extendedchars=false,
  inputencoding=utf8,
}

\usepackage{amsmath}
\usepackage{cite}
\usepackage{tikz}
\usetikzlibrary{positioning,arrows.meta,babel}

\usepackage{enumitem}
\usepackage{multicol}
\usepackage{multirow}
\usepackage{parskip}
\usepackage{placeins}
\usepackage{footmisc}
\usepackage{gensymb}

\bibliographystyle{ieeetr}

\usepackage[bookmarks=true, linkcolor=black, citecolor=black,
            colorlinks=true, urlcolor=black]{hyperref}

% Resolución de referencias cruzadas al documento principal (capítulos,
% secciones, tablas y otros anexos citados desde este anexo).
\usepackage{xr}
\externaldocument{main}

% Encabezado/pie que identifica el documento como anexo del trabajo de grado.
\usepackage{fancyhdr}
\pagestyle{fancy}
\fancyhf{}
\fancyfoot[C]{\footnotesize Anexo del trabajo de grado «FLP Viewer» --- Universidad del Valle, sede Tuluá \quad\textbar\quad \thepage}
\renewcommand{\headrulewidth}{0pt}
\fancypagestyle{plain}{%
  \fancyhf{}%
  \fancyfoot[C]{\footnotesize Anexo del trabajo de grado «FLP Viewer» --- Universidad del Valle, sede Tuluá \quad\textbar\quad \thepage}%
  \renewcommand{\headrulewidth}{0pt}%
}

% El encabezado de capítulo muestra únicamente el título (sin «Capítulo N»),
% ya que cada anexo se publica como documento independiente.
\usepackage{titlesec}
\titleformat{\chapter}[hang]{\normalfont\LARGE\bfseries}{}{0pt}{}
\titlespacing*{\chapter}{0pt}{0pt}{18pt}

\renewcommand{\tablename}{Tabla}
\setcounter{secnumdepth}{3}
\setcounter{tocdepth}{3}
